%
\question[5]
List three anthropogenic effects that can induce alternate stable states in seagrass beds.

\AnswerBox{1\baselineskip}{0.5}{%
	Overfishing, Pollution, and Eutrophication.}


\question[10] A tide pool in the rocky intertidal experiences high turbulence because it is exposed to waves. A nearby tide pool is protected from most waves so it experiences low to moderate turbulence. Both pools have \textit{Enteromorpha} algae, snails, and green crabs. How will the amount of turbulence affect affect the snail's behavior in each pool? How will this affect the abundance of algae present in each pool?

\begin{AnswerPage}{0.5\textheight}%
	High turbulence reduces the ability of the snails to detect the presence of the crabs. (3 points).
	
	\vspace*{\baselineskip}
	
	In a turbulent pool, snails are less likely to detect the crabs so they will forage more often, reducing algal abundance (3 pts).
	
	\vspace*{\baselineskip}
	
	In a quiet pool, snails are more likely to detect crabs so they will forage less, increasing algal abundance. (3 pts)
	
\end{AnswerPage}

%

\question[10]
Would the amount of trophic heat in the ecosystem increase or decrease if climate change causes the greatest temperature difference between immersed and exposed \textit{Pisaster} sea stars.  Explain.

\begin{AnswerPage}{0.6\textheight}%
	Studies showed that feeding rates decreased the most in \textit{Pisaster} when sea stars experienced the greatest thermal stress (greatest temp difference). If feeding rates go down, then trophic heat will increase because they energy they have consumed is lost as metabolic heat instead of being passed to higher trophc levels.	
\end{AnswerPage}

%

\question[10]
\begin{multicols}{2}
	The relative abundance of pico-, nano- and microphytoplankton varies along a gradient from cold temperate to tropical latitudes. The figure at right does not show latitude but does provide support this gradient. Explain how and why, using your knowledge of competitive abilities among the phytoplankton groups and nutrient availability and high and low latitudes.
	
	\columnbreak
	
	\includegraphics[width=0.49\textwidth]{exam2_picoplankton}
	
\end{multicols}	

%

\question[20]
Why is the phytoplankton community in the Gulf of Maine (late spring) composed mostly of microplankton while the community in the middle of the equatorial Atlantic is composed mostly of picoplankton? Your answer must explain how nutrient availability (bottom-up), competition for nutrients, and predation (top-down) contribute to phytoplankton community structure in these two regions. You may use an illustration to aid your explanation.

\begin{AnswerPage}{0.4\textheight}%
	The Gulf of Maine is nutrient rich in the late spring. The equatorial Atlantic is nutrient poor. (2 points)
	
	\vspace*{\baselineskip}
	
	In nutrient poor conditions, picoplankton are better competitors than microplankton so they are abundant. (2 points)
	\vspace*{\baselineskip}
	
	Predators in the equatorial Atlantic are either not an issue or limit the abundance of picoplankton. (2 points)
	
	\vspace*{\baselineskip}
	
	Nutrients never get abundant enough to support large populations of microplankton. (2 points) 
	
	\vspace*{\baselineskip}
	
	In the Gulf, picoplankton abundance becomes limited by predators. That leaves nutrients available for larger plankton sizes. (2 points)
	
	\vspace*{\baselineskip}
	
	Nanoplankton also take up nutrients but their abundance also becomes limited by their predators. (2 points)
	
	\vspace*{\baselineskip}
	
	That leaves sufficient nutrients to support microplankton. (2 points)
	
	\vspace*{\baselineskip}
	
	Microplankton abundance may ultimately get limited by their predators. (2 points)
	
	\vspace*{\baselineskip}
	
	Clarity (2 points)
	
\end{AnswerPage}
